%==============================================================================
\section{Giảm chiều dữ liệu (Dimensionality Reduction)}
\label{sec:dimensionality-reduction}
%==============================================================================

\begin{dinhnghia}[title={Giảm chiều dữ liệu}]
Trong học máy, giảm chiều là quá trình \textbf{giảm số đặc trưng (feature) hoặc biến} trong tập dữ liệu trong khi vẫn giữ được nhiều thông tin gốc nhất có thể. Nói cách khác, đây là cách \textbf{đơn giản hóa} dữ liệu bằng cách giảm độ phức tạp.
\end{dinhnghia}

\begin{dongco}[title={Khi nào cần giảm chiều?}]
Nhu cầu giảm chiều xuất hiện khi dataset có \textbf{rất nhiều} đặc trưng:
\begin{itemize}
  \item Dễ \textbf{overfitting} và làm mô hình phức tạp hơn.
  \item Khó \textbf{trực quan hóa} dữ liệu.
  \item Làm \textbf{chậm} quá trình huấn luyện.
\end{itemize}
\end{dongco}

\subsection{Hai hướng chính}

\begin{phuongphap}[title={Lựa chọn đặc trưng (Feature Selection)}]
Chọn một \textbf{tập con} các đặc trưng gốc theo tiêu chí như mức quan trọng hoặc mức liên quan đến biến mục tiêu.
\end{phuongphap}

\begin{ynghia}[title={Các kỹ thuật lựa chọn đặc trưng phổ biến}]
\begin{itemize}
  \item \textbf{Filter Methods} --- lọc theo thống kê, độc lập mô hình.
  \item \textbf{Wrapper Methods} --- đánh giá hiệu năng mô hình khi thêm/bớt feature.
  \item \textbf{Embedded Methods} --- chọn feature \textbf{gắn trong} quá trình huấn luyện (Lasso, Ridge, cây quyết định, \ldots).
\end{itemize}
\end{ynghia}

\begin{phuongphap}[title={Trích xuất đặc trưng (Feature Extraction)}]
Biến dữ liệu thô thành tập đặc trưng \textbf{có ý nghĩa} hơn cho mô hình: giảm chiều bằng cách chọn, kết hợp hoặc biến đổi feature để tạo tập mới hữu ích hơn.
\end{phuongphap}

\begin{tinhchat}[title={Lợi ích tổng quát}]
Giảm chiều có thể cải thiện \textbf{độ chính xác và tốc độ} mô hình, giảm overfitting và đơn giản hóa trực quan hóa dữ liệu.
\end{tinhchat}
